\begin{table}[t]
\centering
\caption{Light approximate-graph backsolve diagnostics}
\label{tab:approximate-integrability}
\footnotesize
\renewcommand{\arraystretch}{1.08}
\resizebox{0.98\textwidth}{!}{%
\begin{tabular}{>{\raggedright\arraybackslash}p{0.29\textwidth}ccccc}
\toprule
Comparison system & $K$ & Mean abs. residual & p95 abs. residual & Mean envelope width & Max envelope width \\
\midrule
Verified current-cell map & 1 & 0.0000 & 0.0000 & 0.0000 & 0.0000 \\
Verified current-cell map & 5 & 0.0005 & 0.0001 & 0.1801 & 1.6449 \\
Verified current-cell map & 10 & 0.0011 & 0.0006 & 1.3015 & 12.3461 \\
Age-based current matching & 1 & 0.0000 & 0.0000 & 0.0000 & 0.0000 \\
Age-based current matching & 5 & 0.0058 & 0.0401 & 1.2856 & 17.2242 \\
Age-based current matching & 10 & 0.0131 & 0.0780 & 2.4203 & 31.6268 \\
\bottomrule
\end{tabular}
}
\begin{minipage}{0.96\textwidth}
\footnotesize Notes. The table evaluates a light approximate-graph diagnostic in the five-state hump environment at $\rho^{S}=r$. For each comparison system and candidate count $K$, the maintained $K=1$ point summary is treated as a reference potential. The residual is $|\log \eta^{K=1}(x^{\mathrm y})-\log \eta^{K=1}(x^{\mathrm o})-\log(\Qterm(\pi)\WedgeTerm(\pi)/\mathcal{D}^{S}(\pi))|$ on each admissible $K$-nearest edge, weighted by the younger-cell stationary mass. Mean and p95 residuals therefore measure how tightly the wider approximate graph clusters around a single backsolved potential. The last two columns reproduce the maintained envelope widths from the graph-robust exercise so the residual diagnostic can be read alongside the envelope expansion as $K$ grows.
\end{minipage}
\end{table}
